% Rayado college de 7 mm (A4) — Eric Lucero González — https://ericlucero.dev
% Licencia CC BY-NC-SA 4.0. Compilar con pdfLaTeX (dos pasadas: latexmk lo hace solo).
\documentclass[a4paper]{article}
\usepackage[T1]{fontenc}
\usepackage[spanish]{babel}
\usepackage{lmodern}
\usepackage{tikz}
\pagestyle{empty}

% ── Configuración (mm) ─────────────────────────────────────
% El college ruled original es de 9/32 in (7,14 mm); aquí, 7 mm.
\def\anchoPapel{210}
\def\altoPapel{297}
\def\abajo{20}        % espacio libre bajo la última línea
\def\cabecera{32}     % espacio libre arriba (título, fecha)
\def\paso{7}          % distancia entre líneas
\def\margenIzq{30}     % posición de la línea roja de margen

\begin{document}
\begin{tikzpicture}[
    remember picture, overlay, shift={(current page.south west)},
    x=1mm, y=1mm,
    rule lines/.style={cyan!60!gray, very thin},
    margin line/.style={red!60!gray, thin}
  ]
  % Número de renglones que caben entre \abajo y la cabecera
  \pgfmathtruncatemacro{\n}{(\altoPapel - \abajo - \cabecera) / \paso}

  \foreach \i in {0,...,\n} {
    \draw[rule lines] (0, \abajo + \i*\paso) -- (\anchoPapel, \abajo + \i*\paso);
  }

  % Línea de margen
  \draw[margin line] (\margenIzq, 0) -- (\margenIzq, \altoPapel);
\end{tikzpicture}
\end{document}
